% ============================================================
% РАЗДЕЛ «ИСПОЛЬЗОВАНИЕ ПО НАЗНАЧЕНИЮ»
% ============================================================
{\color{unidarkgreen}\section[Использование по назначению]{Использование по назначению}}

% Список кодов для шкафов с ВЧ-функциями (используется в нескольких местах)
\newcommand{\vchcodes}{ЮТКБ.656463.300-011 РЭ, ЮТКБ.656463.300-012 РЭ, ЮТКБ.656463.300-013 РЭ, ЮТКБ.656463.300-014 РЭ, ЮТКБ.656463.300-038 РЭ, ЮТКБ.656463.300-039 РЭ}

% ============================================================
% 1. ЭКСПЛУАТАЦИОННЫЕ ОГРАНИЧЕНИЯ
% ============================================================
\needspace{4\baselineskip}
{\color{unidarkgreen}\subsection{Эксплуатационные ограничения}}

\begin{enumerate}
\item Эксплуатация и обслуживание шкафа должны проводиться в соответствии с НТД РФ — «Правилами технической эксплуатации электрических станций и сетей РФ», «Правил по охране труда при эксплуатации электроустановок», «Правил устройств электроустановок» и требованиями настоящего руководства.
\item Возможность работы шкафа в условиях, отличных от указанных в руководстве, должна согласовываться с предприятием-держателем подлинников конструкторской документации и с предприятием-изготовителем.
\item Условия эксплуатации в части воздействия механических факторов должны соответствовать требованиям настоящего руководства.
\end{enumerate}

% ============================================================
% 2. ПОДГОТОВКА К ЭКСПЛУАТАЦИИ
% ============================================================
{\color{unidarkgreen}\subsection{Подготовка к эксплуатации}}

\subsubsection{Меры безопасности}
\begin{enumerate}
\item Монтаж, обслуживание и эксплуатацию шкафа разрешается производить лицам, прошедшим специальную подготовку, имеющим аттестацию на право выполнения работ (с учётом соблюдения необходимых мер защиты изделий от воздействия статического электричества), хорошо знающим особенности электрической схемы и конструкцию шкафа.
\item Монтаж шкафа и работы на разъёмах терминала, рядах зажимов шкафа и разъёмах устройств следует производить при обесточенном состоянии шкафа. При необходимости проведения проверок с наличием напряжения, должны приниматься дополнительные меры, предотвращающие поражения обслуживающего персонала электрическим током.
\item По требованиям защиты человека от поражения электрическим током шкаф соответствует классу I по ГОСТ 12.2.007.0-75.
\item Запрещается подача какого-либо электроснабжения без заземления. Шкаф перед подачей электроснабжения и во время работы должен быть надежно заземлен.
\end{enumerate}

% Условный подраздел для ВЧ-шкафов
\IfSubStr{\vchcodes}{\devicecode}{%
  \subsubsection{Установка и монтаж ВЧ-приемопередатчика}
  \begin{enumerate}
  \item ВЧ-пост перед установкой должна быть проверена и отрегулирована в соответствии с заводской инструкцией.
  \item Монтаж и подключение ВЧ-поста к шкафу проводится непосредственно на месте эксплуатации шкафа в соответствии со схемой и заводской инструкцией. Установка ВЧ-поста производится в предусмотренное заводом место, согласно чертежа общего вида (ВО).
  \item Подключение цепей связи релейной части защиты с приёмопередатчиком осуществляется проводами, связанными в жгут. На проводах имеется маркировка. Необходимо развязать жгут и соединить провода с выходными зажимами приёмопередатчика.
  \end{enumerate}
}{}%

\subsubsection{Монтаж шкафа}
\begin{enumerate}
\item Шкаф не подвергается консервации смазками и маслами и расконсервации не подлежит.
\item Выполнить подключение шкафа согласно утвержденному проекту в соответствии с указаниями настоящего руководства.
\item Связь шкафа с другими шкафами защит и устройствами производить с помощью кабелей или проводников в соответствии с принятыми проектными техническими решениями.
\item Ввод внешних кабелей осуществляется через гермовводы снизу шкафа. С целью устранения механического усилия натяжения кабеля в шкафу (требование п.2.1.24 ПУЭ), входящие в шкаф кабели вторичных цепей на входе должны быть закреплены зажимом кабельным (к устройству крепления и заземления экранов кабелей) или вводом кабельным с контргайкой типа PG/MG (к панели ввода).
\item Монтаж заземления экранов внешних кабелей необходимо проводить после установки и закрепления шкафа на конструкциях, предусмотренных проектом, и прокладки всех контрольных кабелей.
\item Для заземления экрана кабеля рекомендуется использовать кабельные хомуты из нержавеющей стали или экранирующие зажимы. Кабельные хомуты должны максимально охватывать всю наружную электропроводящую поверхность экрана кабеля и соответствующую этому кабелю перемычку устройства заземления, обеспечивая между ними надёжный электрический контакт.
\item Гермовводы на днище шкафа предназначены для ввода кабелей максимальным диаметром до 20~мм. В случае иных требований, данные требования необходимо указать в картах заказа.
% Условный пункт для ВЧ-шкафов
\IfSubStr{\vchcodes}{\devicecode}{%
  \item Подключение высокочастотной части защиты к высокочастотному каналу связи производить с помощью коаксиального кабеля непосредственно к выходным зажимам приёмопередатчика.
}{}%
\end{enumerate}

\subsubsection{Подготовка шкафа к работе}
\begin{enumerate}
\item Шкаф выпускается с предприятия работоспособным и полностью испытанным. Проверка монтажа и целостности электрических цепей шкафа проведена на предприятии-изготовителе.
\item Положение оперативных переключателей и значения уставок защит шкафа выставить с учётом бланка уставок.
\item Настройки ИЭУ и необходимая оперативная информация доступны на дисплее ИЭУ, при входе в соответствующие пункты меню, с помощью нажатия на соответствующие кнопки управления. Правила работы с ИЭУ подробно описаны в \ManualUnitGeneral
  \ifdefined\ManualUnitDefinite
    и \ManualUnitDefinite.
  \else
    , \ManualUnitDefiniteOne и \ManualUnitDefiniteTwo.
  \fi
\end{enumerate}

% ============================================================
% 3. ИСПОЛЬЗОВАНИЕ ШКАФА
% ============================================================
{\color{unidarkgreen}\subsection{Использование шкафа}}

\begin{enumerate}
\item При вводе шкафа в эксплуатацию необходимо выполнить следующие работы:
  \begin{itemize}
  \item проверку переходного сопротивления заземления шкафа;
  \item проверку сопротивления и прочности изоляции;
  \item проверку коммутационного оборудования в составе шкафа;
  \item задание уставок и конфигурации защит ИЭУ в составе шкафа;
  \item проверку правильности отображения аналоговых величин, поданных от постороннего источника;
  \item проверку уставок измерительных органов (ИО) защит ИЭУ;
  \item проверку воздействия на внешние цепи;
  \item проверку действия на центральную сигнализацию;
  \item проверку взаимодействия шкафа с другими устройствами;
  \item проверку шкафа рабочим током и напряжением.
  \end{itemize}
\item По результатам проверок составляется протокол испытания шкафа при вводе в эксплуатацию (протокол проверок при новом включении).
\end{enumerate}

% ============================================================
% 4. ВОЗМОЖНЫЕ НЕИСПРАВНОСТИ И МЕТОДЫ ИХ УСТРАНЕНИЯ
% ============================================================
{\color{unidarkgreen}\subsection{Возможные неисправности и методы их устранения}}

\begin{enumerate}
\item \textls[-20]{Неисправности могут возникнуть при нарушении условий транспортирования, хранения и эксплуатации.}
\item При включении питания и в процессе работы шкафа могут возникнуть неисправности, обнаруживаемые системой контроля ИЭУ. Описание возможных неисправностей ИЭУ и методов их устранения приведены в руководстве по эксплуатации на ИЭУ.
\item Обслуживающий персонал может самостоятельно провести ремонт или замену коммутационных устройств шкафа, светосигнальной арматуры и т.д.

% Предупреждение — новый абзац (с отступом, если нужно, или без)
\par
\noindent
\begin{tabularx}{\linewidth}{p{2em} X}
  \begin{minipage}[c][3\baselineskip][b]{2em}
    \centering\textcolor{unidarkgreen}{\Huge\faCircleInfo}
  \end{minipage} &
  \hyphenpenalty=0 \textbf{При наличии неисправностей, которые не могут быть исправлены обслуживающим персоналом, и возникшим на стороне непосредственно предприятия-изготовителя, необходимо немедленно уведомить его о возникшей неисправности.}
\end{tabularx}
\end{enumerate}